\section{Reading a formalized research proof with the same method}
We now apply the same method to the research claims described in Chapters~7 and~14. Both appear in OpenAI's public mathematics repository~\cite{oairepo}, which collects preprints together with formal Lean proofs and documents, for each problem, exactly which statements were formalized. A repository listing is not, by itself, a guarantee about every claim in every paper. The right unit of reading is an exact theorem together with its definitions and dependencies, not a title or an impressive final line.

Take Seymour's second-neighborhood conjecture from Chapter~7. The repository's formal statement concerns a nonempty finite oriented graph and a vertex whose second out-neighborhood has at least as many vertices as its first~\cite{oai173}. Chapter~7 supplies the mathematical vocabulary, and Chapter~16 explains how the exclusions in second-neighborhood membership are written in Lean. Reading the formal statement line by line, we find:
\begin{itemize}
\item \code{V} is the type of vertices, and the hypotheses \code{Fintype V}, \code{Nonempty V}, and \code{DecidableEq V} say that there are finitely many vertices, that there is at least one, and that equality of vertices can be decided, which is needed to count finite sets of vertices;
\item the edge relation \code{r} has type \code{V -> V -> Prop}, so \code{r v w} is the proposition $v\to w$;
\item the explicit hypothesis \code{(hr : IsOriented r)} says that \code{r} has no loops and no pair of opposite edges;
\item the conclusion says that there exists a vertex \code{v} whose first neighborhood has at most as many members as its second.
\end{itemize}
Every line corresponds to a part of the informal statement; nothing has been added or left out.

The file that proves the theorem imports three further files, whose names suggest a reduction, a pruning step, and an extremal contradiction. Its proof assumes that a counterexample exists, uses a reduction theorem to obtain a transformed oriented graph with additional properties, and then applies the pruning and extremal results to reach a contradiction~\cite{oaimain}. The reduction produces a relation on pairs of indices from finite sets, with properties named, for instance, positive indegree and strict subset growth~\cite{oaireduction}. Understanding those properties requires opening their definitions; guessing plausible meanings from their names is not enough.

A reading record of this proof therefore has three different kinds of entry. The final statement has been decoded. The logical role of the reduction has been identified. The mathematical proof of the reduction, and the meaning of its additional properties, remain tasks until their source has been read. This book does not reproduce or independently check the repository's argument. It prepares the vocabulary and the method for entering it, without pretending that a short dependency sketch replaces the argument itself.

The repository also contains a configuration file stating precisely what the formal check compares: the file containing the official statement, the file containing the proof, the name of the theorem, and the axioms the proof is permitted to use~\cite{oaiconfig}. The official statement file contains a \code{sorry} in place of a proof, because it only records what must be proved; the check confirms that the separate proof file proves exactly that statement using only the permitted axioms. A matching theorem name in the proof file would not, by itself, be evidence that the two statements agree; the comparison of statements is what matters.

For the circulant Hadamard conjecture of Chapter~14, you can decode the matrix and correlation conditions and prove some necessary restrictions yourself. The repository's description reports a full classification and separates its even-length Barker statement from claims about odd lengths that were not formalized~\cite{oai179}. You can now ask a precise next question: which lemma rules out the candidate orders $4u^2$ with $u>1$ that our elementary argument still permits? That is a mathematically informed question, unlike asking an assistant simply to ``explain the whole repository.''
